%***************************************************************************************************
% preamble/polish.tex - WSZYSTKIE polskie napisy dokumentu w jednym miejscu.
% Ładowany po pakietach (listings, algorithm2e), których nazwy nadpisuje.
%***************************************************************************************************

% Nazwy używane w nagłówkach spisów i podpisach
\newcommand{\abbrevname}{Spis skrótów}

\addto\captionspolish{%
  \renewcommand{\contentsname}{Spis treści}%
  \renewcommand{\listfigurename}{Spis rysunków}%
  \renewcommand{\listtablename}{Spis tabel}%
  \renewcommand{\lstlistlistingname}{Spis listingów}%
  \renewcommand{\chaptername}{Rozdział}%
  \renewcommand{\appendixname}{Załącznik}%
  \renewcommand{\bibname}{Bibliografia}%
  \renewcommand{\refname}{Bibliografia}%
  \renewcommand{\figurename}{Rysunek}%
  \renewcommand{\tablename}{Tabela}%
  \renewcommand{\lstlistingname}{Listing}%
  \renewcommand{\proofname}{Dowód}%
}

% Algorytmy (algorithm2e): nazwa środowiska, nazwa w odwołaniach i tytuł spisu
\SetAlgorithmName{Algorytm}{algorytm}{Spis algorytmów}
\SetAlFnt{\footnotesize}
\SetAlCapFnt{\small}
\SetAlCapNameFnt{\small\bfseries}

% Słowa kluczowe w pseudokodzie (algorithm2e)
\SetKwInOut{Input}{Dane wejściowe}
\SetKwInOut{Output}{Wynik}
\SetKwIF{If}{ElseIf}{Else}{jeżeli}{to}{w p.p. jeżeli}{w p.p.}{}
\SetKwFor{For}{dla}{wykonaj}{}
\SetKwFor{While}{dopóki}{wykonuj}{}
\SetKw{Return}{zwróć}
